\section{Derivación de la constante de estructura fina}
\label{sec:alpha}

La constante de estructura fina recibe en este desarrollo una determinación
aritmética que reúne las tres coordenadas regionales y el registro integral
dodecafásico. La operación combina sus bloques con orientación fijada y
transporta los cocientes de la normalización posicional. El parámetro
resultante expresa la compatibilidad de esa diferencia orientada con
la prolongación de los cilindros. Su valor y los datos de incidencia
utilizados en su construcción permanecerán relacionados durante el
desarrollo excepcional.

El orden interno es importante. Las coordenadas \(\pi\),
\(e\) y \(\varphi\), ya construidas en la sección precedente, son coordenadas
previas del mismo estado aritmético con memoria de trayectoria, y pueden actuar legítimamente como
entradas de un lector posterior. No son constantes convencionales
inyectadas para seleccionar alfa. La generación correlacionada de las cuatro coordenadas comparte origen y
compatibilidad coinductiva, sin que ello imponga evaluar sus cuatro
componentes en un solo paso.

Las dos publicaciones actúan en codominios operatorios distintos. La primera
realiza una clausura entera por división euclídea en base \(1000\), conserva la
cola y produce la sección completa \(\alpha_{\mathrm{HMT}}\). La lectura
analítica tiene dos estratos que no deben confundirse. Su jet de orden nueve,
construido por el resolvente de vacancias \(E_{21/22}\), proporciona una
caracterización finita, aísla una raíz simple y determina una involución
exacta sobre esa fibra. La completación analítica posterior es una
representación correlacionada del mismo estado: calcula todos sus coeficientes
desde las coordenadas regionales y terminales ya publicadas, demuestra la
convergencia uniforme y construye intervalos racionales compatibles con los
cilindros enteros. No se presenta como selector independiente de alfa. La
sección final prueba la igualdad de las dos publicaciones y declara
explícitamente su dependencia común.

La no unicidad de una cola elegida a partir del jet desnudo impide inferir la
completación desde los nueve coeficientes aislados. La regla completa que se
exhibe más abajo elimina esa libertad una vez fijado el estado enriquecido;
no introduce una constante convencional ni altera el orden causal de la
generación dodecafásica de alfa. Estos resultados permanecen en el estado HMT
completo sobre el que se construye el continuo y se formula la demostración
cardinal; sus cifras convencionales no se usan como selectores retroactivos.

\subsection{Primera vía: determinación posicional por clausura entera}
\subsubsection{Dominio y orientación de la clausura entera}
\label{subsec:alpha-dominio}

Sean \(P_j,E_j,\Phi_j\in\{0,\ldots,999\}\) los bloques canónicos de
las tres coordenadas regionales, y prolonguemos el registro dodecafásico por
\[
 K_j=K_{1+((j-1)\bmod12)}.
\]
La orientación del canal es \((1,1,-1)\). La contribución del registro dodecafásico se
sustrae según la clausura dodecafásica ya tipada. Para \(B=1000\), la
operación local distingue el balance firmado
\[
 Z_j:=P_j+E_j-\Phi_j-K_j
\]
de la suma que incorpora el cociente entrante. La normalización es
\begin{equation}
 s_j=Z_j+c_{j+1},\qquad
 A_j=s_j-B\left\lfloor\frac{s_j}{B}\right\rfloor,
 \qquad c_j=\left\lfloor\frac{s_j}{B}\right\rfloor.
 \label{eq:alpha-recurrencia}
\end{equation}
Así, \(A_j\in\{0,\ldots,B-1\}\), incluso cuando \(s_j\) es
negativo. El acarreo que entra desde la derecha pertenece al dominio de
la operación; no puede omitirse porque la tabla se imprima de izquierda
a derecha.

\begin{definition}[Clausura de profundidad finita]
\label{def:alpha-finita}
Para una profundidad \(N\) y una frontera entera \(t\), la aplicación
\(\mathcal C_N\) recibe
\((P_{1:N},E_{1:N},\Phi_{1:N},K_{1:N},c_{N+1}=t)\) y aplica
\eqref{eq:alpha-recurrencia} en el orden \(N,N-1,\ldots,1\). Su salida
es la pareja de palabras
\[
 (A_{1:N},c_{1:N+1}).
\]
\end{definition}

La división euclídea no posee una elección de signo residual: el resto se
fija en \(\{0,\ldots,999\}\) y el cociente queda determinado. Por
ejemplo, \(-431=569-1000\) y \(-1200=800-2\cdot1000\). Los
acarreos negativos de la clausura no son anomalías; son la información
entera necesaria para mantener la igualdad orientada.

\subsubsection{Ventana dodecafásica y condiciones de acarreo}
\label{subsec:alpha-ventana}

Las primeras doce tríadas obtenidas son
\begin{align*}
 P={}&(141,592,653,589,793,238,462,643,383,279,502,884),\\
 E={}&(718,281,828,459,045,235,360,287,471,352,662,497),\\
 \Phi={}&(618,033,988,749,894,848,204,586,834,365,638,117).
\end{align*}
En el canal de clausura, las cinco regiones que determinan la coordenada
\(P\) conservan sus antecedentes distintos. Compartir esta lectura no
las identifica como regiones o como historias.

Con frontera \(c_{13}=0\), la recurrencia produce la tabla siguiente.
Se muestra cada entrada, el acarreo entrante, el total antes de normalizar,
el bloque obtenido y el acarreo saliente. La tabla es un cálculo entero
reproducible, no una lista de bloques seleccionados por su parecido con
un objetivo.

\begin{table}[htbp]
\centering
\small
\setlength{\tabcolsep}{3.5pt}
\begin{tabular}{r|rrrr|rr|rr}
\hline
\(j\)&\(P_j\)&\(E_j\)&\(\Phi_j\)&\(K_j\)&\(c_{j+1}\)&\(s_j\)&\(A_j\)&\(c_j\)\\
\hline
1&141&718&618&234&0&7&007&0\\
2&592&281&033&543&0&297&297&0\\
3&653&828&988&140&\(-1\)&352&352&0\\
4&589&459&749&729&\(-1\)&\(-431\)&569&\(-1\)\\
5&793&045&894&659&\(-2\)&\(-717\)&283&\(-1\)\\
6&238&235&848&824&\(-1\)&\(-1200\)&800&\(-2\)\\
7&462&360&204&621&0&\(-3\)&997&\(-1\)\\
8&643&287&586&058&\(-1\)&285&285&0\\
9&383&471&834&914&\(-1\)&\(-895\)&105&\(-1\)\\
10&279&352&365&794&0&\(-528\)&472&\(-1\)\\
11&502&662&638&146&0&380&380&0\\
12&884&497&117&601&0&663&663&0\\
\hline
\end{tabular}
\caption{Clausura inicial con frontera derecha nula. La dependencia se propaga de la última fila a la primera.}
\label{tab:alpha-ventana}
\end{table}

La palabra de acarreos es
\[
 (c_1,\ldots,c_{13})=(0,0,0,-1,-1,-2,-1,0,-1,-1,0,0,0),
\]
y la lectura de la ventana es el racional
\begin{equation}
 \alpha_{12}^{(0)}
 =\frac{7297352569283800997285105472380663}{10^{36}}
 =0.007297352569283800997285105472380663.
 \label{eq:alpha-ventana-cero}
\end{equation}
El superíndice recuerda la frontera elegida. Esta ventana es exacta en su
dominio finito; no se la identifica sin más con el truncamiento de la
sección infinita. La diferencia entre ambas lecturas se hará visible al
transportar el acarreo de la cola remota.

\subsubsection{Telescopado y reversibilidad en una fibra de frontera}
\label{subsec:alpha-telescopia}

La forma local de la igualdad es
\begin{equation}
 P_j+E_j-\Phi_j-K_j-A_j=Bc_j-c_{j+1}.
 \label{eq:alpha-identidad-local}
\end{equation}
El término de acarreo es un enlace entre posiciones contiguas. Si se lo
elimina coordenada a coordenada, se obtiene una igualdad falsa. Sólo una
evaluación que conserve sus extremos permite telescoparlo.

\begin{proposition}[Identidad entera de profundidad arbitraria]
\label{prop:alpha-telescopia}
Para \(\iota_N(z)=\sum_{j=1}^{N}z_jB^{N-j}\), toda ejecución de
\(\mathcal C_N\) satisface
\begin{equation}
 \iota_N(P)+\iota_N(E)-\iota_N(\Phi)-\iota_N(K)-\iota_N(A)
 =B^Nc_1-c_{N+1}.
 \label{eq:alpha-telescopia}
\end{equation}
\end{proposition}

\begin{proof}
Multiplicar \eqref{eq:alpha-identidad-local} por \(B^{N-j}\) y sumar
produce en el lado derecho
\[
 \sum_{j=1}^{N}(Bc_j-c_{j+1})B^{N-j}=B^Nc_1-c_{N+1}.
\]
Cada acarreo interior aparece dos veces con signo opuesto; sólo sobreviven
los dos extremos.
\end{proof}

Para la tabla~\ref{tab:alpha-ventana}, ambos extremos son cero. Si
\(p_{12},e_{12},f_{12}\) son las tres sumas fraccionarias truncadas,
\[
 \alpha_{12}^{(0)}=p_{12}+e_{12}-f_{12}-\kappa_{36}.
\]
La constante del registro dodecafásico que aparece es \(\kappa_{36}\), no
\(\kappa_{\mathrm{per}}\). El dominio de esta identidad contiene
truncamientos y una frontera finita; su forma no autoriza a borrar ambas
condiciones al pasar a la sección completa.

La inversión local es también exacta:
\[
 K_j=P_j+E_j-\Phi_j+c_{j+1}-A_j-Bc_j.
\]
Dados \(P,E,\Phi\), los bloques \(A\) y los acarreos completos,
recupera el registro dodecafásico. La concatenación recupera además
\[
 \iota_N(K)=\iota_N(P)+\iota_N(E)-\iota_N(\Phi)-\iota_N(A)
             -B^Nc_1+c_{N+1}.
\]
Esta reversibilidad es una propiedad de la fibra con datos y fronteras
fijados. No debe convertirse en una nueva receta causal que empiece por
alfa y fabrique después el registro. La dirección de construcción sigue
siendo registro firmado, registro dodecafásico, clausura y salida.

\subsubsection{Prolongación y normalización canónica de la sección completa}
\label{subsec:alpha-completa}

La prolongación conserva el registro dodecafásico periódico y recibe las secuencias
regionales completas. Sean
\[
 p=\sum_{j\geq1}P_jB^{-j},\qquad
 e_0=\sum_{j\geq1}E_jB^{-j},\qquad
 f=\sum_{j\geq1}\Phi_jB^{-j}.
\]
Cada serie es una realización del carácter interno correspondiente.
Las partes enteras ya obtenidas son 3, 2 y 1, de modo que
\(p=\pi-3\), \(e_0=e-2\) y
\(f=\varphi-1\).

\begin{proposition}[Existencia y unicidad de la normalización completa]
\label{prop:alpha-normalizacion}
La serie
\begin{equation}
 y=p+e_0-f-\kappa_{\mathrm{per}}
   =\sum_{j\geq1}(P_j+E_j-\Phi_j-K_j)B^{-j}
 \label{eq:alpha-serie}
\end{equation}
converge absolutamente. En el dominio de los bloques obtenidos,
\(0<y<10^{-2}\). Existe una única expansión canónica \((A_j)\) de
\(y\), sin cola finalmente constante 999, y una única sucesión entera
acotada \((c_j)\) que satisface \eqref{eq:alpha-identidad-local} con
\(c_1=0\). Estos acarreos pertenecen al conjunto invariante
\(\mathcal C=\{-2,-1,0,1,2\}\).
\end{proposition}

\begin{proof}
Los coeficientes de la serie están acotados por \(2(B-1)\) en módulo,
de modo que domina una serie geométrica. Si
\(Z_j=P_j+E_j-\Phi_j-K_j\) y
\(R_j(Z)=\sum_{r\geq0}Z_{j+r}B^{-r-1}\), cada cola es la diferencia
de dos sumas de dos colas canónicas y pertenece a \((-2,2)\).
El primer coeficiente es
\(Z_1=141+718-618-234=7\); por tanto
\(y=(7+R_2(Z))/1000\in(0.005,0.009)\).

Sea \((A_j)\) la expansión canónica de \(y\). Definamos
\[
 c_j=R_j(Z)-R_j(A).
\]
Al multiplicar la igualdad de las dos series por \(B^{j-1}\) y separar
sus primeros \(j-1\) términos se obtiene
\[
 c_j=\sum_{r=1}^{j-1}(A_r-Z_r)B^{j-1-r}\in\mathbb Z,
\]
y \(c_1=0\). Además, \(R_j(A)\in[0,1)\), luego
\(-3<c_j<2\), lo que incluye todos los acarreos en \(\mathcal C\).
Las identidades de desplazamiento de las colas dan
\(Bc_j=Z_j-A_j+c_{j+1}\), que es la recurrencia requerida.

Recíprocamente, cualquier solución acotada con salida canónica y
\(c_1=0\) telescopa en el límite a la misma serie \(y\). La expansión
canónica es única y, una vez fijada, la fórmula de las colas determina
todos los acarreos. Finalmente, si se aplica una ejecución finita con
\(c_{j+1}\in\mathcal C\), se tiene \(-2000\leq s_j\leq2000\),
por lo que su cociente euclídeo vuelve a pertenecer a \(\mathcal C\).
\end{proof}

La proposición explicita la normalización del lector posterior: no
reemplaza la generación de las tres secuencias por una serie convencional
introducida como dato. Una vez producidas esas secuencias por el estado
coinductivo, su combinación y la frontera de la normalización quedan
determinadas sin consultar alfa.

El procedimiento finito conservado en el corpus propaga las cinco
fronteras de \(\mathcal C\) y conserva solamente el prefijo común a
todas ellas. Si dos profundidades han coalescido antes de una frontera,
la recurrencia determinista de derecha a izquierda da los mismos bloques
a su izquierda. Ésa es la compatibilidad con el truncamiento. No se
escoge una frontera porque dé cifras favorables, ni se deduce la
compatibilidad a toda profundidad de haber observado una corrida larga.

La comparación inicial muestra por qué esta precaución importa. En la
prolongación estabilizada se obtiene \(c_{13}^{\infty}=-1\), mientras
que la ventana aislada imponía \(c_{13}=0\). Así,
\[
 884+497-117-601-1=662
\]
en el último bloque de la primera vuelta. Las primeras once tríadas no
cambian, pero la prolongación comienza
\[
 (007,297,352,569,283,800,997,285,105,472,380,662,999,\ldots).
\]
Una aparición del bloque 999 no viola la convención canónica: lo
excluido es una cola que sea finalmente constante 999. La ventana con
frontera cero termina en 663; el truncamiento de la sección completa
termina en 662. Ambas afirmaciones describen operaciones diferentes y
exactas. No deben fusionarse en una sola notación de proyección.

\subsubsection{Definición operativa y ecuación aditiva completa}
\label{subsec:alpha-definicion}

\begin{definition}[Coordenada alfa de la clausura]
\label{def:alpha-operativa}
La vía entera determina la coordenada
\begin{equation}
 \alpha_A=\sum_{j\geq1}A_jB^{-j}
 =p+e_0-f-\kappa_{\mathrm{per}},
 \label{eq:alpha-via-a}
\end{equation}
donde \((A_j,c_j)\) es la normalización canónica compatible de
\eqref{eq:alpha-recurrencia}. En la realización escalar de las
coordenadas HMT,
\begin{equation}
 \boxed{\alpha_A=
 \pi+e-\varphi
 -4-\kappa_{\mathrm{per}}.}
 \label{eq:alpha-identidad-completa}
\end{equation}
\end{definition}

El entero cuatro no se calibra: es
\(3+2-1\), combinación de las partes enteras de las tres coordenadas.
La base mil procede de la carta cúbica y el período doce del registro dodecafásico. Los
signos son los del canal orientado y del cierre, y la frontera remota es
la que impone la sección compatible. Éstos son los orígenes efectivos de
los coeficientes de la ecuación.

En este nivel puede describirse alfa como la coordenada escalar del
defecto de clausura entre las tres realizaciones regionales y el retorno
racional dodecafásico, calculado con acarreo compatible. El término
\emph{defecto} designa una diferencia algebraica determinada, no una
pérdida arbitraria ni una imperfección del sistema. Esta descripción no
suprime la genealogía: sin las regiones, el registro firmado, la
orientación y la ley de acarreo, la igualdad escalar no sería una
exposición suficiente de la construcción.

También deben mantenerse separados tres objetos que la notación de alfa
puede acompañar en lectores posteriores: \(A\) es una palabra de
bloques; el soporte negativo de preacarreo se obtiene del signo de los
totales antes de normalizar; y un vector radial excepcional pertenece a
un codominio de incidencia. Ninguno de ellos es intercambiable con el
valor real \(\alpha_A\).

\paragraph{Balance firmado y suma con cociente entrante.}
\label{ampl-alfa-z-s-aclaracion}
La división posicional consta de dos etapas consecutivas. Con
\(K_j^{\mathrm{per}}=K_{1+((j-1)\bmod12)}\), que designa la
prolongación periódica denotada también por \(K_j\) en
\eqref{eq:alpha-recurrencia}, el balance anterior al transporte del
cociente es
\[
 Z_j:=P_j+E_j-\Phi_j-K_j^{\mathrm{per}}.
\]
La variable $s_j$ de \eqref{eq:alpha-recurrencia} designa la suma que ya
incluye la contribución entrante, por lo que
\[
 s_j=Z_j+c_{j+1},\qquad
 A_j=Z_j+c_{j+1}-1000c_j.
\]
La configuración incidencial
$H^-_\alpha$ utiliza los signos del balance $Z_j$ de la ventana inicial;
la determinación de las cifras utiliza además los cocientes de enlace
$c_{j+1}$ y $c_j$. Esta distinción conserva conjuntamente el soporte
firmado y la normalización de la expansión.

% Frontera de truncamiento de la vía dodecafásica completa y delimitación
% exacta del resolvente analítico finito.
\paragraph{Dependencia exacta respecto de la frontera de truncamiento.}
\label{ampl-alfa-frontera-exacta}
Las coordenadas regionales y el registro dodecafásico proceden de la
generación común APP--TRIT--TPK. Una vez producidas, su composición
posicional posee una dependencia explícita respecto de la profundidad
de lectura. Sean $B=1000$,
\[
 Z_j=P_j+E_j-\Phi_j-K_j^{\mathrm{per}},\qquad
 y_N=\sum_{j=1}^{N}Z_jB^{-j},\qquad
 R_{N+1}(Z)=\sum_{r\geq0}Z_{N+1+r}B^{-r-1}.
\]
La serie completa de \eqref{eq:alpha-serie} satisface exactamente
\[
 \alpha_A=y_N+B^{-N}R_{N+1}(Z),\qquad
 -2<R_{N+1}(Z)<2.
\]
Esta identidad reúne la contribución de la ventana y la de su
prolongación, ambas expresadas en las mismas coordenadas.

\begin{proposition}[Transporte de la frontera y cilindro canónico]
\label{ampl-alfa-prop-frontera}
Sea $A_j^{(N,t)}$ la salida de la división posicional finita con
$c_{N+1}=t\in\mathcal C=\{-2,-1,0,1,2\}$, y sea
$a_N^{(t)}=\sum_{j=1}^{N}A_j^{(N,t)}B^{-j}$. Para los registros del
dominio considerado, se cumple
\begin{equation}
 \alpha_A-a_N^{(t)}
 =B^{-N}\bigl(R_{N+1}(Z)-t\bigr),
 \qquad
 |\alpha_A-a_N^{(t)}|<4B^{-N}.
 \label{ampl-alfa-eq-frontera}
\end{equation}
El cociente de la prolongación completa es
$c_{N+1}^{\infty}=\lfloor R_{N+1}(Z)\rfloor$. Con esta frontera,
\[
 0\leq\alpha_A-a_N^{(c_{N+1}^{\infty})}<B^{-N},
\]
de modo que la ventana obtenida es el prefijo canónico de orden $N$.
\end{proposition}
\begin{proof}
El telescopado de \eqref{eq:alpha-identidad-local} da
$a_N^{(t)}=y_N+tB^{-N}-c_1$. Como $Z_1=7$, cada estado entrante
pertenece a $\mathcal C$ y el primer dividendo queda entre $5$ y $9$,
se tiene $c_1=0$. Sustituir esta igualdad en la descomposición de la
serie prueba \eqref{ampl-alfa-eq-frontera}. La cota resulta de
$R_{N+1}(Z)\in(-2,2)$ y $t\in\mathcal C$.

Para la normalización canónica completa,
$c_{N+1}^{\infty}=R_{N+1}(Z)-R_{N+1}(A)$ y
$R_{N+1}(A)\in[0,1)$. La integralidad del cociente implica
$c_{N+1}^{\infty}=\lfloor R_{N+1}(Z)\rfloor$. La fórmula de error
se convierte entonces en
$\alpha_A-a_N^{(c_{N+1}^{\infty})}=B^{-N}R_{N+1}(A)$, que prueba
la pertenencia al cilindro canónico.
\end{proof}

La proposición identifica el significado de la frontera derecha: es la
parte entera de la cola firmada expresada en la escala de la posición
siguiente. El cálculo sobre las cinco fronteras conserva esta información
como una familia finita de posibilidades hasta que los bloques a la
izquierda se estabilizan. La contracción cuantitativa procede de la
profundidad y del transporte de la cola; su objeto es la normalización de
las coordenadas ya generadas.

\paragraph{Dominio graduado y estatuto de la representación analítica.}
\label{ampl-alfa-dominio-resolvente}
La composición analítica exhibida en
\S\ref{subsec:alpha-resolvente} actúa sobre
$\operatorname{span}\{e_7,e_8,e_9\}$ y satisface $Ne_9=0$. Su
evaluación produce exactamente $T_{7:9}$; el grado $9$ es la última
componente de ese dominio. La identidad $N^3=0$ explica por qué el
resolvente se expresa mediante tres sumandos. El resolvente local no determina
por sí mismo coeficientes de grado superior. La
sección~\ref{subsec:alpha-limite-dos-vias} demuestra esa no unicidad del jet
desnudo y separa dos operaciones: la publicación dodecafásica produce la
coordenada \(\alpha_{\mathrm{HMT}}\) desde los registros del estado; después,
la representación analítica consume la misma precoordenada y genera, por una
recurrencia explícita, todos los coeficientes de su imagen analítica. No se
consulta una raíz convencional, ni se promueve la representación analítica a
generador independiente. Los cuadrados de compatibilidad y los cilindros
anidados prueban la igualdad de ambas publicaciones a cada profundidad. La
nilpotencia local, la publicación dodecafásica y la representación analítica
desempeñan así funciones distintas: la primera certifica el resultado finito
de orden nueve; la segunda determina \(\alpha_{\mathrm{HMT}}\); la tercera
representa esa misma coordenada a toda profundidad y verifica su compatibilidad.


\subsubsection{Determinación constructiva y función del registro dodecafásico}
\label{alpha-ampl-dependencias}

La definición anterior permite precisar qué significa determinar alfa desde
la estructura. El dominio de la operación contiene cuatro registros de
procedencia conocida: las tres secuencias regionales y el registro integral
de fase. La normalización incorpora, además, una condición de frontera
compatible con la prolongación. Una vez fijados esos datos, la división
euclídea determina cada resto y su cociente. La proposición
\ref{prop:alpha-normalizacion} extiende esa determinación al sistema
completo. La independencia respecto de un valor físico objetivo es, por
tanto, una propiedad del orden de construcción: el valor que se compara al
final no interviene en la selección de ninguno de los argumentos de esta
aplicación.

Esta independencia causal debe distinguirse de la independencia algebraica
entre números. La primera se establece identificando el dominio, las
operaciones y la procedencia de los coeficientes; la segunda sería una
afirmación sobre relaciones polinómicas. El resultado utilizado aquí es el
primero. La construcción determina una coordenada mediante registros
producidos por APP--TRIT--TPK y permite comprobar posteriormente las
relaciones que satisface. Ese orden hace posible utilizar una coordenada ya
generada en una operación ulterior sin convertirla en un parámetro externo.

El registro $K$ tiene una función más precisa que la de una corrección
numérica añadida a una fórmula. Sus doce posiciones proceden de la
transformación reversible del registro firmado. La ecuación
$U=\mathcal H_{12}K$ conserva la posibilidad de volver a ese registro,
mientras las diferencias transversales y la suma total permiten otra
reconstrucción. En la evaluación periódica, el orden de las posiciones
determina el peso posicional de cada componente. El mismo conjunto de
componentes, dispuesto en otro orden, produce en general otra coordenada
racional. Por ello el origen de fase y la orientación forman parte de la
definición del número que interviene en \eqref{eq:alpha-identidad-completa}.

La normalización de alfa compone esa información con las regiones. Su
identidad local distingue dos operaciones sucesivas. Primero se forma el
balance orientado de bloques; después se representa ese balance mediante
un resto canónico y un cociente que enlaza posiciones contiguas. La primera
operación conserva el signo del balance y permite la extracción posterior
de un soporte. La segunda permite formar los prefijos de la coordenada
real. Ambas lecturas permanecen relacionadas porque el estado conserva los
datos anteriores y posteriores a la normalización. El soporte firmado y la
expansión posicional son, así, realizaciones diferentes de una composición
aritmética explícita.

La identidad telescópica expresa el contenido global de esta composición.
Los cocientes interiores se cancelan por pares al evaluar un prefijo,
mientras sus extremos permanecen en la igualdad. Esa cancelación es una
propiedad de la suma ponderada; no elimina los cocientes del estado ni
permite suprimirlos antes de efectuar la suma. La frontera derecha de una
ventana recibe la contribución de la continuación. Por esa razón, la
evaluación de una ventana aislada y la restricción de una sección
prolongada son operaciones distintas. El cambio del bloque final $663$ a
$662$, demostrado anteriormente, ofrece una realización explícita de esa
diferencia sin alterar los bloques ya estabilizados a su izquierda.

En este sentido, la bidireccionalidad tiene una formulación aritmética
concreta. La generación proporciona prolongaciones compatibles; la
normalización transmite la información entera de su frontera desde las
posiciones de menor peso hacia las de mayor peso. La ecuación conserva la
contribución de ambas direcciones mediante sus índices y sus condiciones
de contorno. La reversibilidad en la fibra fijada se comprueba reconstruyendo
$K$ a partir de los demás registros y de los cocientes completos. Esa
inversión certifica la consistencia de la operación original; la dirección
genealógica continúa siendo la que produce primero el registro y después
la coordenada de alfa.

La publicación analítica correlacionada tiene un contenido complementario.
En ella, el objeto inmediato es un balance de vacancias y la coordenada ya
producida por el estado aparece representada como raíz en un dominio
aislante. La publicación dodecafásica proporciona una normalización canónica
de registros; la representación analítica estudia la anulación de un
operador sobre esa misma coordenada. La exposición conjunta conserva la
diferencia de codominios sin afirmar independencia causal. Sus relaciones
con el estado común se expresan mediante las aplicaciones de refinamiento y
la comparación de sus publicaciones completas, tal como se desarrolla en la
subsección correspondiente. El acuerdo de una aproximación finita constituye
un control localizado; la igualdad a toda profundidad requiere los cilindros
anidados y la recurrencia coinductiva demostrados allí.

Finalmente, esta organización explica la posición de alfa en el artículo.
Su ecuación aparece después de la generación regional y de la construcción
de $K$, porque utiliza ambos resultados. La sección electrónica utiliza
después esa coordenada como un factor de su composición. La prolongación
excepcional recibe los registros y sus soportes, donde se conserva la
información de incidencia. Estas dos continuaciones retoman una misma
procedencia y desarrollan funciones matemáticas distintas: evaluación de
una sección central y construcción de una configuración de frontera.
El desarrollo de cada una permite seguir la estructura a través de la
ecuación, en lugar de limitar la comparación a su valor escalar.


\subsection{Lectura analítica: testigo finito y representación correlacionada}
\subsubsection{Invariantes del transporte y resolvente graduado}
\label{subsec:alpha-via-b}

El testigo analítico finito recibe invariantes del estado orientado, no la
salida \(\alpha_A\). Su truncamiento de orden seis es
\begin{equation}
 P_6(x)=2\pi x-\frac74x^2
 +\frac{x^3}{2\pi}
 +\frac{x^4}{20}-\frac{2x^5}{21}-\frac{x^6}{46},
 \label{eq:alpha-p6}
\end{equation}
y la condición de vacancia utiliza
\[
 \delta=\log_{10}\frac{10}{9}.
\]
La razón \(10/9\) compara la carta decimal con las nueve posiciones
APP de la construcción nonádica. Su logaritmo pertenece a esta lectura
posterior de compactificación; no es un valor objetivo de alfa.

La genealogía de los coeficientes debe conservarse completa, pero con su
fuerza lógica exacta. Los desarrollos recientes reúnen una firma marcada
del estado. La matriz de sincronización y el censo asociado son
\[
 D_{\mathrm{syn}}=\begin{pmatrix}85&112\\41&52\end{pmatrix},
 \qquad N_9=43.
\]
De su determinante se extrae
\[
 a=-\frac{\det D_{\mathrm{syn}}}{N_9}
   =\frac{172}{43}=4.
\]
El reloj inducido sobre los huecos cortos tiene espectro secundario
\(\{6,7\}\); su extremo superior da \(b=7\). La cardinalidad
trítica \(m=3\) determina
\[
 t_*=mb=21,\qquad k_*=mb-1=20.
\]
La lectura de vacancia en la carta mil da
\[
 v_-=\lfloor1000\delta\rfloor=45,\qquad
 v_+=\lceil1000\delta\rceil=46.
\]
Estos enteros son antecedentes tipados de la evaluación del truncamiento. No se
obtienen de su raíz ni de una aproximación metrológica a ella.

La rama inferior coincide además con el determinante del bloque APP
\[
 B_c=\begin{pmatrix}7&2\\2&7\end{pmatrix},\qquad\det B_c=45.
\]
El rango transversal del registro del capítulo anterior da
\(r_{\mathrm{dod}}=11\), y la semicorona de fase proporciona
\(f_{\mathrm{ph}}=8\binom62=120\). En particular,
\(45r_{\mathrm{dod}}=495\).

Con \(\omega=2\pi\), la firma
\[
 \mathcal I_9=(\omega,m,a,b,k_*,t_*,v_-,v_+,f_{\mathrm{ph}},r_{\mathrm{dod}})
\]
alimenta el lector graduado
\begin{align}
 \mathfrak E_9(\mathcal I_9)(x)
 ={}&\omega x-\frac ba x^2+\omega^{-1}x^3
 +\frac{x^4}{k_*}-\frac{2x^5}{t_*}-\frac{x^6}{v_+}\notag\\
 &-\frac{x^7}{f_{\mathrm{ph}}}+\frac{x^8}{v_-}
 +\frac{2x^9}{v_-r_{\mathrm{dod}}}.
 \label{eq:alpha-firma-coeficientes}
\end{align}
La sustitución recupera los coeficientes de \(P_6\) y de su
prolongación de orden nueve. El par \(45/46\) atraviesa la separación
entre ambos: la rama superior ocupa el grado seis y la inferior reaparece
en los grados ocho y nueve.

La factorización de \eqref{eq:alpha-firma-coeficientes} está formulada
para un estado que conserva origen, orden de canales y orientación. La
naturalidad respecto de morfismos que preservan esos datos no equivale a
la unicidad de todo lector imaginable después de olvidarlos. En el dominio
marcado, la fórmula anterior fija cada coeficiente por
\((\omega,m,a,b,k_*,t_*,v_-,v_+,f_{\mathrm{ph}},r_{\mathrm{dod}})\);
por tanto, dos lectores que preserven la firma y apliquen la misma
evaluación graduada coinciden término a término hasta el grado nueve. Esta
es la unicidad que se utiliza a continuación y no exige una remisión
documental adicional.

\subsubsection{Resolvente nilpotente en los grados \(7\), \(8\) y \(9\)}
\label{subsec:alpha-resolvente}

La primera prolongación admite una construcción operatoria finita
especialmente explícita. Sobre el espacio con base \(e_7,e_8,e_9\),
sea
\[
 Ne_7=-\frac83e_8,\qquad Ne_8=\frac2{11}e_9,\qquad Ne_9=0,
 \qquad v_7=-\frac1{120}e_7.
\]
El símbolo \(N\) designa aquí un operador nilpotente, no una profundidad
de truncamiento. Su dominio y graduación están fijados antes de evaluar
el polinomio en un número.

\begin{lemma}[Prolongación nilpotente exacta]
\label{lem:alpha-resolvente}
Se cumple \(N^3=0\) y
\[
 (I-N)^{-1}v_7
 =-\frac{e_7}{120}+\frac{e_8}{45}+\frac{2e_9}{495}.
\]
La evaluación \(e_n\mapsto x^n\) produce
\begin{equation}
 T_{7:9}(x)=-\frac{x^7}{120}+\frac{x^8}{45}+\frac{2x^9}{495}.
 \label{eq:alpha-t79}
\end{equation}
\end{lemma}

\begin{proof}
El operador aumenta el grado y se anula en el último vector de la base,
por lo que su cubo es cero. La inversa es el resolvente finito
\(I+N+N^2\). Además,
\[
 Nv_7=\frac1{45}e_8,\qquad N^2v_7=\frac2{495}e_9.
\]
Sumar los tres términos y evaluar da la fórmula.
\end{proof}

Aquí los denominadores 120, 45 y 495 no son ajustes decimales: proceden
de la semicorona, del determinante APP y de su producto por el rango
dodecafásico. Las transferencias \(-8/3\) y \(2/11\), junto con el
signo y el grado de la semilla, pertenecen al lector tipado. Conservar
sólo la simetría del ciclo y sustituir este operador por otro operador
simétrico no conserva necesariamente los coeficientes.

\subsubsection{Raíz de multiplicidad uno del truncamiento de orden nueve}
\label{subsec:alpha-jet}

Definamos
\[
 P_9=P_6+T_{7:9},\qquad E_{21/22}^{(9)}(x)=P_9(x)-\delta,
 \qquad I_\alpha=(0,10^{-2}).
\]
Este operador finito tiene un resultado propio de existencia y unicidad.

Para aislar su raíz mediante aritmética racional, sean
\[
 A_n(t)=\sum_{j=0}^{n}\frac{(-1)^jt^{2j+1}}{2j+1},\qquad
 L_n(t)=2\sum_{j=0}^{n}\frac{t^{2j+1}}{2j+1}\quad(0<t<1).
\]
El criterio alternante y la suma geométrica de la cola dan
\begin{align}
 A_{2n+1}(t)&<\arctan t<A_{2n}(t),
 \label{eq:alpha-cotas-arctan}\\
 L_n(t)&<\log\frac{1+t}{1-t}
 <L_n(t)+\frac{2t^{2n+3}}{(2n+3)(1-t^2)}.
 \label{eq:alpha-cotas-log}
\end{align}
Aplicamos \eqref{eq:alpha-cotas-arctan} a la identidad de Machin
\(\pi=16\arctan(1/5)-4\arctan(1/239)\), y
\eqref{eq:alpha-cotas-log} a
\[
 \log(10/9)=\log\frac{1+1/19}{1-1/19},\qquad
 \log 10=\log(10/9)+2\log\frac{1+1/2}{1-1/2}.
\]
Escribimos
\[
 \delta:=\frac{\log(10/9)}{\log 10}.
\]
Con los órdenes de truncamiento \(80,20,40,120\), aplicados
respectivamente a las cuatro series anteriores, se obtienen intervalos
racionales para \(\pi\) y \(\delta\), ambos de anchura menor que
\(10^{-70}\). No se consulta ningún valor de alfa en estas cotas.

\begin{proposition}[Raíz de multiplicidad uno del truncamiento]
\label{prop:alpha-raiz-jet}
El operador \(E_{21/22}^{(9)}\) tiene una única raíz
\(\xi_9\) en \(I_\alpha\), y esa raíz es simple. Además posee el
aislamiento racional
\begin{equation}
 \frac{729735256928380099}{10^{20}}
 <\xi_9<\frac{729735256928380100}{10^{20}}.
 \label{eq:alpha-intervalo-jet}
\end{equation}
\end{proposition}

\begin{proof}
Las cotas de las coordenadas internas
\(3.14<\pi<3.15\) y
\(0.045<\delta<0.046\) son suficientes para la existencia. En cero,
\(E_{21/22}^{(9)}(0)=-\delta<0\). En \(10^{-2}\), la contribución
lineal menos los términos negativos es mayor que \(0.0626\), luego el
valor del operador es positivo.

Para \(0\leq x\leq10^{-2}\), al descartar los términos positivos de
la derivada se obtiene
\[
 (E_{21/22}^{(9)})'(x)
 \geq2\pi-\frac72\cdot10^{-2}
 -\frac{10}{21}\cdot10^{-8}-\frac6{46}\cdot10^{-10}
 -\frac7{120}\cdot10^{-12}>6.24.
\]
La continuidad da existencia; la cota derivativa da unicidad y
simplicidad. Para el extremo izquierdo \(l\) y el extremo derecho \(u\)
de \eqref{eq:alpha-intervalo-jet}, la sustitución racional de las cotas
\eqref{eq:alpha-cotas-arctan}--\eqref{eq:alpha-cotas-log} en el polinomio
da, con redondeo dirigido,
\[
 -4.6\,10^{-20}<E_{21/22}^{(9)}(l)<-4.5\,10^{-20}<0,
 \qquad
 0<1.6\,10^{-20}<E_{21/22}^{(9)}(u)<1.8\,10^{-20}.
\]
La monotonía ya probada sitúa, por tanto, la única raíz entre ambos
extremos. El intervalo se verifica después de construir el operador y no
interviene en la selección de sus coeficientes.
\end{proof}

El resultado concierne a \(\xi_9\). No autoriza a escribir
\(\xi_9=\alpha_A\) como una igualdad exacta de números completos.
La raíz de un truncamiento puede compartir un prefijo con la raíz del
operador prolongado y diferir en una posición posterior. La simplicidad
permite controlar el movimiento de la raíz bajo perturbaciones; no
anula esas perturbaciones.

\subsection{Compatibilidad entre las publicaciones dodecafásica y analítica de alfa}
\label{subsec:alpha-limite-dos-vias}

La compatibilidad se establece a toda profundidad. La publicación
dodecafásica reversible y la representación analítica de la
coordenada de alfa
\(E_{21/22}\) son dos publicaciones internas con codominios distintos,
producidas a partir del mismo estado enriquecido APP--TRIT--TPK. El jet
\(E_{21/22}^{(9)}\) es un testigo finito exacto, pero no determina por sí solo
la cola. La completación que se construye después no se deduce del jet
desnudo: recibe la coordenada compartida del estado y publica su imagen
analítica mediante una recurrencia explícita. Por ello no constituye una
selección causal independiente; constituye una representación analítica
tipada del mismo punto generado.

Esta separación conserva el orden causal
\[
 \mathrm{APP}\longrightarrow\mathrm{TRIT}\longrightarrow\mathrm{TPK}
 \longrightarrow\mathcal X_{\alpha}^{\mathrm{enr}}
 \longrightarrow (P,E,\Phi,U_{12}^{\mathrm{sgn}},K)
 \longrightarrow\alpha_{\mathrm{HMT}}.
\]
Las coordenadas \(P,E,\Phi\) y el registro \(K\) son salidas anteriores del
mismo estado enriquecido. Su uso en la clausura de alfa es, por tanto, una
composición interna posterior y no una inyección de constantes
convencionales.

Para \(r\in6\mathbb N\), \(r\ge36\), sea
\(\widehat{\mathcal X}^{\mathrm{enr},\alpha}_r\) el sector de estados
truncados que prolongan la semilla interna de alfa fijada en \(R_{36}\). Se
escribe
\[
 x_r^\alpha=(\widetilde x_r,\chi_\alpha,
              \varepsilon_r^\alpha,N_r^\alpha),
\]
donde \(\chi_\alpha\) es el carácter interno del sector, no el valor escalar
publicado. Cada estado conserva residuo, cociente, acarreo, hoja, orientación,
fase, ruta, cilindro, frontera, firma, registro, memoria, supervivencia,
procedencia y profundidad. Si
\[
 d_r^\alpha=\operatorname{Dig}_{729}(\varepsilon_r^\alpha),\qquad
 \varepsilon_{r+6}^\alpha=729\varepsilon_r^\alpha-d_r^\alpha,
 \qquad N_{r+6}^\alpha=729N_r^\alpha+d_r^\alpha,
\]
la transición efectiva es
\[
 \mathcal U_r^\alpha=
 \operatorname{Upd}_r^\alpha\circ
 \operatorname{Tra}_r^\alpha\circ
 \operatorname{Sel}_r^\alpha,
 \qquad
 \operatorname{Ext}^{\alpha,\to}_{r+6,r}
 =\operatorname{graph}(\mathcal U_r^\alpha).
\]
La funcionalidad de estas extensiones, con la semilla \(R_{36}\) fijada,
produce una única sección compatible: la fase retorna mientras memoria y
profundidad avanzan. El dominio global de ambas determinaciones es
\begin{equation}
 \mathcal X_\alpha^{\mathrm{enr}}:=\widehat\Omega_\alpha
 :=\varprojlim_{r\ge36}
 \bigl(\widehat{\mathcal X}^{\mathrm{enr},\alpha}_r,
       \operatorname{Ext}^{\alpha,\to}_{r+6,r}\bigr).
 \label{eq:alpha-dominio-enriquecido-completo}
\end{equation}

\subsubsection{La clausura dodecafásica como sección completa}

Para \(N\geq1\), pongamos
\[
 \mathcal W_N:=\{0,\ldots,999\}^{N},
 \qquad
 \mathsf T_{\alpha,N}(\mathsf s):=(A_1,\ldots,A_N),
 \qquad \mathsf s\in\mathcal X_{\alpha}^{\mathrm{enr}},
\]
donde los bloques \(A_j\) se obtienen mediante
\eqref{eq:alpha-recurrencia} con la frontera compatible
\(c_{N+1}^{\infty}=\lfloor R_{N+1}(Z)\rfloor\). Si \(N'\geq N\), sea
\(\rho_{N',N}:\mathcal W_{N'}\to\mathcal W_N\) la restricción a los
primeros \(N\) bloques.

\begin{theorem}[Salida dodecafásica completa y compatibilidad proyectiva]
\label{thm:alpha-salida-dodecafasica-completa}
La familia \((\mathsf T_{\alpha,N})_{N\geq1}\) es compatible:
\begin{equation}
 \rho_{N',N}\circ\mathsf T_{\alpha,N'}
 =\mathsf T_{\alpha,N}
 \qquad(N'\geq N).
 \label{eq:alpha-via-a-proyectiva-exacta}
\end{equation}
Sus cilindros canónicos
\begin{equation}
 C_N^A(\mathsf s):=
 \left[
  \sum_{j=1}^{N}A_j1000^{-j},
  \sum_{j=1}^{N}A_j1000^{-j}+1000^{-N}
 \right)\cap I_\alpha
 \label{eq:alpha-cilindros-via-a-exactos}
\end{equation}
son anidados, tienen diámetro a lo sumo \(1000^{-N}\) y satisfacen
\begin{equation}
 \bigcap_{N\geq1}C_N^A(\mathsf s)
 =\{\alpha_A(\mathsf s)\},
 \qquad
 \alpha_A
 =\pi_{\mathrm{HMT}}+e_{\mathrm{HMT}}-\varphi_{\mathrm{HMT}}
  -4-\kappa_{\mathrm{per}}.
 \label{eq:alpha-salida-a-completa}
\end{equation}
En consecuencia, la primera vía publica una única sección infinita, que
denotamos \(\alpha_{\mathrm{HMT}}:=\alpha_A\).
\end{theorem}

\begin{proof}
La proposición~\ref{prop:alpha-normalizacion} proporciona la expansión
canónica única y la sucesión acotada única de acarreos. La identidad local
\eqref{eq:alpha-identidad-local} y su telescopado
\eqref{eq:alpha-telescopia} conservan la frontera remota; la proposición
\ref{ampl-alfa-prop-frontera} prueba que restringir una ejecución prolongada
recupera exactamente el prefijo anterior. Esto da
\eqref{eq:alpha-via-a-proyectiva-exacta}. La pertenencia de la suma completa
a cada cilindro se sigue de la identidad de colas, y
\(\operatorname{diam}C_N^A\leq1000^{-N}\to0\); la intersección contiene un
solo punto. Finalmente, \eqref{eq:alpha-identidad-completa} identifica ese
punto con la expresión de \eqref{eq:alpha-salida-a-completa}. Todos sus
coeficientes y signos se fijan antes de la publicación de alfa.
\end{proof}

La ventana \(\alpha_{12}^{(0)}\), cerrada con \(c_{13}=0\), sigue siendo un
cálculo finito exacto; no debe confundirse con la restricción de la sección
completa, cuya frontera compatible satisface \(c_{13}^{\infty}=-1\). En
particular,
\begin{equation}
 \rho_{36,12}(\alpha_{A,36})=\alpha_{A,12}
 \label{eq:alpha-rho-36-12-proyectiva}
\end{equation}
es una igualdad de palabras compatibles, no una igualdad entre palabras de
longitudes distintas ni una identificación de la ventana aislada con el
límite.

\subsubsection{Caracterización analítica exacta de orden nueve}

Sea
\[
 \mathbb K:=\mathbb Q\!\left(\pi_{\mathrm{HMT}},
                  \log_{10}\frac{10}{9}\right),
 \qquad
 E_9(x):=E_{21/22}^{(9)}(x)=P_9(x)-\delta .
\]
Las ecuaciones \eqref{eq:alpha-firma-coeficientes} y
\eqref{eq:alpha-t79} construyen \(P_9=P_6+T_{7:9}\) desde la firma del
estado y desde el resolvente nilpotente; la
proposición~\ref{prop:alpha-raiz-jet} demuestra que \(E_9\) posee una única
raíz simple \(\xi_9\in I_\alpha\), aislada por el intervalo racional
\eqref{eq:alpha-intervalo-jet}.

La bidireccionalidad finita puede escribirse sin introducir una cola
analítica. Definamos
\begin{align}
 B_9(x):={}&\delta+\frac74x^2-\frac{x^4}{20}
 +\frac{2x^5}{21}+\frac{x^6}{46}+\frac{x^7}{120}
 -\frac{x^8}{45}-\frac{2x^9}{495},
 \label{eq:alpha-b9-finito}\\
 q_x(T):={}&xT^2-B_9(x)T+x^3,\qquad x\in I_\alpha=(0,10^{-2}),
 \label{eq:alpha-q-involucion-finita}\\
 J_x:{\mathbb R}^{\times}\longrightarrow{\mathbb R}^{\times},
 \qquad J_x(T):={}&\frac{x^2}{T}.
\end{align}

\begin{proposition}[Equivalencia cuadrática e involución finitas]
\label{prop:alpha-involucion-finita-exacta}
Para todo \(x\in I_\alpha\) y todo \(T\in\mathbb R^\times\),
\begin{align}
 \frac{q_x(2\pi_{\mathrm{HMT}})}{2\pi_{\mathrm{HMT}}}
 &=P_9(x)-\delta=E_9(x),
 \label{eq:alpha-q-equivale-p9}\\
 J_x(J_x(T))&=T,
 \qquad
 T^2q_x(J_x(T))=x^2q_x(T).
 \label{eq:alpha-j-involucion-exacta}
\end{align}
Sea \(Z_x=\{T\in\mathbb R^\times:q_x(T)=0\}\). Por consiguiente,
\(J_x|_{Z_x}:Z_x\to Z_x\) es involutiva. Cuando \(E_9(x)=0\),
intercambia las raíces \(2\pi_{\mathrm{HMT}}\) y
\(x^2/(2\pi_{\mathrm{HMT}})\). Esta es una involución exacta sobre la fibra
escalar del jet de orden nueve; no actúa sobre el estado enriquecido.
\end{proposition}

\begin{proof}
Dividir \(q_x(T)\) por \(T\) y sustituir
\(T=2\pi_{\mathrm{HMT}}\) reproduce término a término la definición de
\(P_9(x)-\delta\). La primera identidad de
\eqref{eq:alpha-j-involucion-exacta} es inmediata. Para la segunda,
\[
 T^2q_x(x^2/T)
 =x^5-B_9(x)x^2T+x^3T^2=x^2q_x(T).
\]
Así, \(J_x\) preserva el conjunto de raíces y, por la fórmula de Viète,
intercambia las dos raíces cuyo producto es \(x^2\).
\end{proof}

El testigo decimal finito conserva entonces su fuerza exacta:
\begin{equation}
 \operatorname{Tr}_9(\xi_9^{-1})
 =\operatorname{Tr}_9\!\left((\alpha_{12}^{(0)})^{-1}\right)
 =137.035999084.
 \label{eq:alpha-testigo-finito-nueve}
\end{equation}
La igualdad \eqref{eq:alpha-testigo-finito-nueve} certifica la
caracterización finita y la involución precedente. No convierte la raíz de
un polinomio truncado en la raíz de toda prolongación compatible.

\subsubsection{El jet no determina por sí solo una prolongación analítica}

Sea
\[
 j_9:\mathbb K[[x]]\longrightarrow \mathbb K[x]/(x^{10})
\]
el truncamiento, y para \(h\in\mathbb K[x]\) pongamos
\begin{equation}
 E_h(x):=E_9(x)+x^{10}h(x).
 \label{eq:alpha-familia-prolongaciones-mismo-jet}
\end{equation}
Todos estos polinomios satisfacen \(j_9(E_h)=j_9(E_9)\).

\begin{theorem}[No unicidad de la raíz a partir del jet]
\label{thm:alpha-no-unicidad-prolongacion-jet}
No existe una aplicación
\[
 \mathfrak r_9:\mathbb K[x]/(x^{10})\longrightarrow I_\alpha
\]
que asigne, a partir del jet solamente, la raíz positiva de toda
prolongación polinómica de \(E_9\) que tenga una única raíz en
\(I_\alpha\). En particular, \(E_0\) y \(E_{1/729}\) poseen el mismo jet y
raíces positivas distintas.
\end{theorem}

\begin{proof}
La cota demostrada en la prueba de
\ref{prop:alpha-raiz-jet} da \(E_9'(x)>6.24\) en
\([0,10^{-2}]\). Para \(t\geq0\), escribamos
\(E_t(x)=E_9(x)+tx^{10}\). Entonces
\[
 E_t'(x)=E_9'(x)+10tx^9>6.24,
\]
mientras \(E_t(0)=-\delta<0\) y
\(E_t(10^{-2})>E_9(10^{-2})>0\). Cada \(E_t\) posee, por tanto, una raíz
positiva única \(\eta(t)\in I_\alpha\). Para \(t=1/729\),
\[
 E_{1/729}(\xi_9)=\frac{\xi_9^{10}}{729}>0;
\]
como \(E_{1/729}\) es estrictamente creciente y cambia de signo una sola
vez, su raíz satisface \(\eta(1/729)<\xi_9=\eta(0)\). Sin embargo,
\(j_9(E_{1/729})=j_9(E_0)\). Una aplicación que factorizase la raíz por
\(j_9\) tendría que asignar el mismo punto a ambos polinomios, contradicción.
\end{proof}

La variación no es sólo existencial. El teorema de la función implícita se
aplica porque \(\partial_xE_t(\eta(t))>6.24\), y proporciona
\begin{equation}
 \eta'(t)=-
 \frac{\eta(t)^{10}}
 {E_9'(\eta(t))+10t\eta(t)^9}<0.
 \label{eq:alpha-variacion-raiz-prolongacion}
\end{equation}
Así queda probado que los coeficientes posteriores al grado nueve modifican
la raíz aun cuando dejan invariante todo el jet construido hasta ese grado.
Esta conclusión no niega la cola producida por el estado enriquecido: prueba
únicamente que la raíz no factoriza por el truncamiento desnudo
\(j_9\). Escoger una cola \(h\) a partir de una raíz convencional externa
sería circular. La completación que sigue no consulta esa raíz ni una tabla
metrológica: evalúa directamente las coordenadas regionales y terminales del
estado. Como esa evaluación coincide algebraicamente con la publicación
entera, el resultado es una representación correlacionada, no una afirmación
de independencia entre dos selectores.

% Representación analítica correlacionada de alfa: acción, convergencia e
% intervalos.
\subsubsection{Representación analítica de la coordenada de alfa}
\label{subsec:alpha-lector-completo-rev08}

El control del jet anterior obliga a exhibir la cola: no basta nombrar un
lector completo. Sea \(\mathsf s\in\mathcal X_\alpha^{\mathrm{enr}}\). Antes
de la normalización decimal, el mismo estado ya publica la precoordenada
\begin{equation}
 a(\mathsf s):=p(\mathsf s)+e_0(\mathsf s)-f(\mathsf s)
                 -\kappa_{\mathrm{per}}(\mathsf s)\in I_\alpha,
 \label{eq:alpha-precoordenada-comun}
\end{equation}
con \(p=\pi_{\mathrm{HMT}}-3\),
\(e_0=e_{\mathrm{HMT}}-2\) y
\(f=\varphi_{\mathrm{HMT}}-1\). Equivalentemente,
\(a=\pi_{\mathrm{HMT}}+e_{\mathrm{HMT}}-
\varphi_{\mathrm{HMT}}-4-\kappa_{\mathrm{per}}\). No se introduce aquí
\(\alpha\) convencional: \(p,e_0,f\) son los tres caracteres regionales ya
generados y \(\kappa_{\mathrm{per}}\) es la lectura racional del registro
terminal. La clausura entera aplica a
\eqref{eq:alpha-precoordenada-comun} la normalización de
\eqref{eq:alpha-recurrencia}; la carta que sigue publica sobre la misma
precoordenada una representación analítica compatible. Su grafo de
dependencia recibe \((p,e_0,f,\kappa_{\mathrm{per}})\) directamente: no
contiene una arista \(\alpha_A\to\lambda\), aunque la igualdad demostrada más
abajo permita identificar después ambos nombres de publicación.

Escribamos \(q=1/729\), sea \(F_{\mathsf s}\) el cuerpo ordenado generado por
las coordenadas de \(\mathsf s\), y pongamos
\(E_{9,\mathsf s}(X)=E_{21/22}^{(9)}(\mathsf s;X)\). Definimos el coeficiente
de enlace
\begin{equation}
 \lambda(\mathsf s):=
 -\frac{E_{9,\mathsf s}(a(\mathsf s))
         \bigl(1-q\,a(\mathsf s)^3\bigr)}{a(\mathsf s)^{10}}.
 \label{eq:alpha-lambda-estado}
\end{equation}
Todos sus argumentos pertenecen al estado antes del reconocimiento
metrológico. En particular, \(\lambda\) no se obtiene ajustando una cadena
decimal objetivo. La ecuación~\eqref{eq:alpha-lambda-estado} debe leerse con
su tipo exacto: una vez que la primera publicación del estado ha producido
\(a(\mathsf s)\), fija la única prolongación de la clase geométrica que se
define a continuación. No constituye una segunda selección independiente de
\(a(\mathsf s)\), sino una representación analítica correlacionada del mismo
estado enriquecido.

En efecto, para \(\mu\in F_{\mathsf s}\) pongamos
\[
 E_{\mathsf s,\mu}(X)
 :=E_{9,\mathsf s}(X)+\mu\frac{X^{10}}{1-qX^3}.
\]
Como \(0<a(\mathsf s)<10^{-2}\) y \(qa(\mathsf s)^3<1\), se tiene
\begin{equation}
 E_{\mathsf s,\mu}(a(\mathsf s))=0
 \quad\Longleftrightarrow\quad
 \mu=\lambda(\mathsf s).
 \label{eq:alpha-unicidad-lambda-clase-geometrica}
\end{equation}
Así, la semilla de coeficientes no contiene libertad una vez fijados el estado
y la clase de prolongaciones \(q\)-geométricas. Esta unicidad es relativa a la
clase declarada; no afirma que el jet de orden nueve seleccione por sí solo una
cola entre todas las series analíticas posibles.

La acción local sobre bloques de tres coeficientes es
\begin{equation}
 \mathcal C_\alpha:F_{\mathsf s}^{3}\longrightarrow F_{\mathsf s}^{3},
 \qquad \mathcal C_\alpha(u,v,w)=q(u,v,w),
 \qquad c^{(0)}=(\lambda(\mathsf s),0,0).
 \label{eq:alpha-accion-coeficientes}
\end{equation}
Por tanto, para todo \(n\geq0\),
\begin{equation}
 b_{10+3n}=q^n\lambda(\mathsf s),\qquad
 b_{11+3n}=b_{12+3n}=0.
 \label{eq:alpha-recurrencia-coeficientes}
\end{equation}
Ésta es la regla que calcula cada coeficiente posterior al grado nueve. No
queda un parámetro libre en cada prolongación.

Para \(M\geq0\), sea
\begin{equation}
 E_{\mathsf s,M}(X):=E_{9,\mathsf s}(X)
 +\lambda(\mathsf s)\sum_{n=0}^{M}q^nX^{10+3n}.
 \label{eq:alpha-truncamientos-completos}
\end{equation}
La función completa es
\begin{equation}
 E_{\mathsf s,\infty}(X)
 :=E_{9,\mathsf s}(X)
 +\lambda(\mathsf s)\frac{X^{10}}{1-X^3/729}.
 \label{eq:alpha-funcion-completa-explicita}
\end{equation}

Para formular el sistema proyectivo sin tipos implícitos, fijamos
\(d_M=12+3M\) y definimos la fibración de jets
\begin{equation}
 \mathfrak J_M:=
 \coprod_{\mathsf s\in\mathcal X_\alpha^{\mathrm{enr}}}
 \{\mathsf s\}\times F_{\mathsf s}[X]/(X^{d_M+1}).
 \label{eq:alpha-espacio-jets-rev08}
\end{equation}
Si \(M'\ge M\), la restricción es
\begin{equation}
 \tau_{M',M}(\mathsf s,[P]_{d_{M'}})
 :=(\mathsf s,[P]_{d_M}).
 \label{eq:alpha-truncamiento-jets-rev08}
\end{equation}

\begin{proposition}[Naturalidad de la recurrencia analítica]
\label{prop:alpha-naturalidad-recurrencia-explicita}
Sea
\[
 \Gamma_M(\mathsf s)
 =\bigl(\mathsf s,[E_{\mathsf s,M}]_{d_M}\bigr).
\]
Si \(M'\geq M\), el truncamiento \(\tau_{M',M}\) satisface
\begin{equation}
 \tau_{M',M}\circ\Gamma_{M'}=\Gamma_M,
 \qquad
 \Gamma_{E21}(\mathsf s)=\varprojlim_M\Gamma_M(\mathsf s).
 \label{eq:alpha-naturalidad-explicita}
\end{equation}
La proyección sobre los tres grados nuevos es precisamente
\(\mathcal C_\alpha^{M+1}(c^{(0)})\).
\end{proposition}

\begin{proof}
La ecuación~\eqref{eq:alpha-recurrencia-coeficientes} fija el bloque de
grados \(10+3n,11+3n,12+3n\). Al pasar de \(M\) a \(M+1\) sólo se añade
el bloque siguiente; truncarlo recupera literalmente el polinomio anterior.
La compatibilidad para \(M'\geq M\) se obtiene por composición. Así, el
límite inverso no es una sección escogida del jet desnudo, sino la órbita
única de \(c^{(0)}\) bajo \(\mathcal C_\alpha\).
\end{proof}

\subsubsection{Convergencia uniforme, raíz simple e intervalos racionales}
\label{subsec:alpha-cotas-completas-rev08}

Los cilindros regionales y el registro periódico proporcionan cotas mediante
aritmética de intervalos racionales. Para cada coordenada
\[
 c\in\{\pi_{\mathrm{HMT}},e_{\mathrm{HMT}},\varphi_{\mathrm{HMT}}\},
\]
sea \(P_c\) su prefijo nonádico certificado de mil cifras. El dato utilizado
no es el racional \(P_c\) como si fuese el valor completo, sino el cilindro
semiabierto y su cierre exterior
\[
 \mathcal C_c^{(1000)}=[P_c,P_c+\varepsilon),
 \qquad
 \overline{\mathcal C}_c^{(1000)}=[P_c,P_c+\varepsilon],
 \qquad \varepsilon=10^{-1000}.
\]
La aritmética dirigida opera conservadoramente con esos cierres. Si
\(P_\pi,P_e,P_\varphi\) denotan los tres prefijos, se obtiene
\begin{equation}
\begin{aligned}
 a_-&:=P_\pi+P_e-(P_\varphi+\varepsilon)-4-
       \kappa_{\mathrm{per}},\\
 a_+&:=(P_\pi+\varepsilon)+(P_e+\varepsilon)-P_\varphi-4-
       \kappa_{\mathrm{per}},\\
 A_{1000}&:=[a_-,a_+],
 \qquad a(\mathsf s)\in[a_-,a_+]=A_{1000},
 \qquad \operatorname{diam}A_{1000}=3\varepsilon.
\end{aligned}
\label{eq:alpha-propagacion-colas-rev08}
\end{equation}
La extensión intervalar natural de
\eqref{eq:alpha-lambda-estado} transporta después \(A_{1000}\), el cilindro
de \(\pi_{\mathrm{HMT}}\) y la cota racional dirigida de \(\delta\) a un
intervalo cerrado \(\Lambda_{1000}\ni\lambda(\mathsf s)\); no se congela
ninguna de las tres colas. En particular,
\begin{equation}
\begin{aligned}
 \frac{7297352569283800997285105472380662}{10^{36}}
 &<a(\mathsf s)<
 \frac{7297352569283800997285105472380663}{10^{36}},\\
 -8.711\mathbin{\times}10^{-6}
 &<\lambda(\mathsf s)<-8.709\mathbin{\times}10^{-6}.
\end{aligned}
 \label{eq:alpha-cotas-precoordenada-lambda}
\end{equation}
Estas desigualdades usan prefijos certificados con su cota racional de
cola; no reciben una tabla de alfa. El verificador del paquete las reproduce
desde las tres salidas nonádicas y desde la carta terminal.

Sea
\[
 r:=\frac{10^{-6}}{729}<1.372\mathbin{\times}10^{-9}.
\]
Para \(0\leq x\leq10^{-2}\), la cola posterior a \(M\) satisface
\begin{equation}
 \left|E_{\mathsf s,\infty}(x)-E_{\mathsf s,M}(x)\right|
 \leq
 8.711\mathbin{\times}10^{-26}\,\frac{r^{M+1}}{1-r}.
 \label{eq:alpha-cota-cola-uniforme}
\end{equation}
Además,
\begin{equation}
 \left|\frac{\mathrm d}{\mathrm dx}
  \left(\lambda\frac{x^{10}}{1-x^3/729}\right)\right|
 \leq 8.711\mathbin{\times}10^{-6}
 \left(\frac{10\,10^{-18}}{1-r}
 +\frac{3\,10^{-24}/729}{(1-r)^2}\right)
 <8.712\mathbin{\times}10^{-23}.
 \label{eq:alpha-cota-derivada-cola}
\end{equation}
La cola de las derivadas posterior al bloque \(M\) admite, además, la
estimación dependiente de la profundidad
\begin{equation}
 \sup_{0\le x\le10^{-2}}
 \left|E'_{\mathsf s,\infty}(x)-E'_{\mathsf s,M}(x)\right|
 \le 8.711\mathbin{\times}10^{-24}\,r^{M+1}
 \left(\frac{13+3M}{1-r}+\frac{3r}{(1-r)^2}\right)
 \xrightarrow[M\to\infty]{}0.
 \label{eq:alpha-cota-resto-derivadas-rev08}
\end{equation}

\begin{theorem}[Función límite y raíz simple única de la carta correlacionada]
\label{thm:alpha-raiz-completa-explicita}
La sucesión \(E_{\mathsf s,M}\) converge uniformemente, junto con sus
derivadas, a \(E_{\mathsf s,\infty}\) en
\(\overline I_\alpha=[0,10^{-2}]\). Se cumple
\begin{equation}
 E_{\mathsf s,\infty}(a(\mathsf s))=0,
 \qquad
 \inf_{x\in\overline I_\alpha}
 E'_{\mathsf s,\infty}(x)>6.239.
 \label{eq:alpha-raiz-y-derivada-completas}
\end{equation}
Por consiguiente, \(a(\mathsf s)\) es la raíz simple única de la función
completa en \(I_\alpha\).
\end{theorem}

\begin{proof}
La razón de dos términos consecutivos de la cola está acotada por \(r<1\),
lo que da \eqref{eq:alpha-cota-cola-uniforme}; la derivación término a
término y la suma de las dos series geométricas dan
\eqref{eq:alpha-cota-derivada-cola} y
\eqref{eq:alpha-cota-resto-derivadas-rev08}. La prueba del jet establece
\(E'_{9,\mathsf s}>6.24\) en \(\overline I_\alpha\); por tanto, la
perturbación completa deja la derivada por encima de \(6.239\).
Finalmente, al sustituir \eqref{eq:alpha-lambda-estado} en
\eqref{eq:alpha-funcion-completa-explicita}, los dos sumandos evaluados en
\(a(\mathsf s)\) se cancelan exactamente. La función es estrictamente
creciente; su única raíz es, pues, simple y coincide con esa precoordenada.
\end{proof}

Para hacer efectiva la identificación a toda profundidad, las evaluaciones
dirigidas de los extremos certifican primero
\[
 E_{\mathsf s,\infty}(0)<0<
 E_{\mathsf s,\infty}(10^{-2}),
 \qquad
 E'_{\mathsf s,\infty}\geq\mu:=6239/1000>c:=6
 \quad\hbox{en }[0,10^{-2}].
\]
Por continuidad existe una raíz en el intervalo y, por la segunda desigualdad,
es única. Sólo después de fijar esta existencia, la monotonía y la invariante
de que la horquilla contiene la raíz se aplica el rescate por el teorema del
valor medio. Partimos de \(D_0=[0,10^{-2}]\). Si
\(D_j=[r_j,s_j]\), sea \(m_j=(r_j+s_j)/2\), y sea
\([F_j^-,F_j^+]\) una inclusión racional dirigida de
\(E_{\mathsf s,\infty}(m_j)\), obtenida propagando los cilindros de
\(\pi_{\mathrm{HMT}},e_{\mathrm{HMT}},\varphi_{\mathrm{HMT}}\),
\(\delta\) y \(\lambda\), con
\begin{equation}
 0\leq F_j^+-F_j^-\leq \frac{c}{4}\,(s_j-r_j).
 \label{eq:alpha-anchura-evaluacion-rescate-rev08}
\end{equation}
Cuando \(F_j^+<0\) se conserva la mitad derecha;
cuando \(F_j^->0\), la mitad izquierda. Si
\(F_j^-\leq0\leq F_j^+\), no se intenta decidir si el punto medio es
exactamente la raíz. Sea \(w_j=s_j-r_j\). La cota uniforme
\(E'_{\mathsf s,\infty}\geq c=6\) y el teorema del valor medio implican que
cualquier raíz contenida en \(D_j\) pertenece a
\begin{equation}
 D_{j+1}=D_j\cap
 \left[m_j-\frac{w_j}{4},
             m_j+\frac{w_j}{4}\right].
 \label{eq:alpha-rescate-derivada-rev08}
\end{equation}
En efecto, sea \(\xi_j\) la raíz y sea
\(y_j=E_{\mathsf s,\infty}(m_j)\in[F_j^-,F_j^+]\). Como el recinto contiene
cero y satisface \eqref{eq:alpha-anchura-evaluacion-rescate-rev08}, se tiene
\(|y_j|\leq cw_j/4\). El teorema del valor medio proporciona
\(|m_j-\xi_j|\leq |y_j|/c\leq w_j/4\), lo que prueba
\eqref{eq:alpha-rescate-derivada-rev08}. El argumento incluye el caso
\(y_j=0\): la raíz está entonces en el punto medio, pero el algoritmo no
necesita reconocer esa igualdad para conservarla.

Si el recinto disponible no satisface
\eqref{eq:alpha-anchura-evaluacion-rescate-rev08}, se refinan primero los
cilindros fuente y se repite la evaluación; no se declara un signo ni se
acepta una contracción insuficiente. En cualquiera de las tres ramas el nuevo
intervalo conserva la raíz y tiene diámetro a lo sumo \(w_j/2\). Para cada
profundidad objetivo finita, el refinamiento dirigido permite satisfacer la
cota de anchura sin comparar de forma no certificada un número real con cero.
La monotonía demuestra inductivamente
\begin{equation}
 a(\mathsf s)\in D_{j+1}\subseteq D_j,\qquad
 \operatorname{diam}D_{j+1}\leq
 \tfrac12\operatorname{diam}D_j,
 \qquad \operatorname{diam}D_j\longrightarrow0,
 \label{eq:alpha-biseccion-racional}
\end{equation}
donde la desigualdad es una cota de contracción y no una igualdad impuesta.
El certificado ejecutable aplica
esta regla a veinticuatro niveles en base mil y comprueba que el cierre
exterior completo \(A_{1000}\), no sólo su extremo inferior, queda dentro de
cada cilindro publicado. Sus pruebas negativas rechazan un recinto demasiado
ancho, un recinto desordenado, un punto que no sea el medio, un dominio sin
cambio de signo y una horquilla degenerada. La prolongación coinductiva permite
reemplazar mil por cualquier profundidad requerida.

Concretamente, para \(S_N=1000^N\) el verificador calcula
\begin{equation}
 j_N^-:=\lfloor S_Na_-\rfloor,
 \qquad j_N^+:=\lfloor S_Na_+\rfloor
 \label{eq:alpha-indices-cilindro-dirigidos-rev08}
\end{equation}
y sólo publica el nivel si \(j_N^-=j_N^+\). Esta igualdad, que incluye el
extremo superior del cierre exterior, prueba
\(A_{1000}\subset[j_N^-/S_N,(j_N^-+1)/S_N)\) y evita atribuir a la celda
izquierda un extremo situado exactamente sobre su frontera derecha.

Sea \(C_N^A(\mathsf s)\) el cilindro entero de
\eqref{eq:alpha-cilindros-via-a-exactos}. Elegimos recursivamente
\(m(N)>m(N-1)\) de modo que
\(\operatorname{diam}D_{m(N)}<1000^{-N}/4\), y definimos
\begin{equation}
I_{B,N}:=D_{m(N)}\cap C_N^A(\mathsf s).
 \label{eq:alpha-intervalos-b-completos}
\end{equation}
Definimos
\(\ell_N:=\inf I_{B,N}\) y \(u_N:=\sup I_{B,N}\). Ambos son racionales;
el intervalo es no vacío porque los dos factores contienen
\(a(\mathsf s)\). Puesto que el cilindro es semiabierto, \(u_N\) puede no
pertenecer a \(I_{B,N}\). La evaluación en ese extremo se entiende mediante
la extensión continua de \(E_{\mathsf s,\infty}\) al cierre exterior. Así,
\begin{equation}
 I_{B,N+1}\subseteq I_{B,N}\subseteq C_N^A(\mathsf s),
 \qquad \operatorname{diam}I_{B,N}\longrightarrow0,
 \qquad E_{\mathsf s,\infty}(\ell_N)\leq0\leq
 E_{\mathsf s,\infty}(u_N).
 \label{eq:alpha-inclusion-signos-completa}
\end{equation}
La cota~\eqref{eq:alpha-cota-cola-uniforme} permite escoger un
truncamiento que conserve el signo en cada extremo donde la función completa
no se anula. Si un extremo coincide con la raíz, se utiliza la identidad de
la función completa y la pertenencia ya demostrada; no se atribuye ese mismo
signo a sus truncamientos. En efecto, el resto exacto es
\[
 E_{\mathsf s,\infty}(x)-E_{\mathsf s,M}(x)
 =\lambda(\mathsf s)x^{10}
   \frac{(qx^3)^{M+1}}{1-qx^3}.
\]
En el dominio positivo considerado, donde \(\lambda(\mathsf s)<0\), se
obtiene en la raíz
\[
 E_{\mathsf s,M}(a(\mathsf s))
 =-\lambda(\mathsf s)a(\mathsf s)^{10}
   \frac{(qa(\mathsf s)^3)^{M+1}}{1-qa(\mathsf s)^3}>0
 \qquad(M<\infty).
\]
Por ello, la certificación de los signos débiles en
\eqref{eq:alpha-inclusion-signos-completa} se refiere a
\(E_{\mathsf s,\infty}\). Su evaluación utiliza la función completa o el
truncamiento acompañado del resto con signo. Esto cubre tanto el ínfimo
perteneciente como el supremo eventualmente excluido del cilindro.

\begin{proposition}[Cuadrados de compatibilidad y cilindros comunes]
\label{prop:alpha-cuadrados-compatibilidad}
Las restricciones de palabras y de publicaciones analíticas conmutan:
\[
\begin{array}{ccc}
 \mathcal X_\alpha^{\mathrm{enr}}&
 \xrightarrow{\ \mathsf T_{\alpha,N'}\ }&\mathcal W_{N'}\\[1mm]
 \Big\Vert&&\Big\downarrow\rho_{N',N}\\[1mm]
 \mathcal X_\alpha^{\mathrm{enr}}&
 \xrightarrow{\ \mathsf T_{\alpha,N}\ }&\mathcal W_N
\end{array}
\qquad
\begin{array}{ccc}
 \mathcal X_\alpha^{\mathrm{enr}}&
 \xrightarrow{\ \Gamma_{M'}\ }&\mathfrak J_{M'}\\[1mm]
 \Big\Vert&&\Big\downarrow\tau_{M',M}\\[1mm]
 \mathcal X_\alpha^{\mathrm{enr}}&
 \xrightarrow{\ \Gamma_M\ }&\mathfrak J_M .
\end{array}
\]
La familia \(I_{B,N}\) proporciona la conexión efectiva entre ambos
cuadrados y satisface materialmente
\(I_{B,N}\subseteq C_N^A\) para todo \(N\).
\end{proposition}

\begin{proof}
El primer cuadrado es la compatibilidad de la normalización entera; el segundo
es la proposición~\ref{prop:alpha-naturalidad-recurrencia-explicita}. La
inclusión se sigue de la definición~\eqref{eq:alpha-intervalos-b-completos},
y la no vacuidad y los signos se siguen del
teorema~\ref{thm:alpha-raiz-completa-explicita}. Ningún paso identifica la
raíz de un truncamiento con la de una cola arbitraria.
\end{proof}

Definimos ahora, sin dejar implícito el lector de prefijos,
\begin{equation}
 \alpha_B(\mathsf s):=
 \operatorname*{arg\,zero}_{x\in I_\alpha}E_{\mathsf s,\infty}(x),
 \qquad
 \operatorname{pref}_N(x):=
 \frac{\lfloor1000^Nx\rfloor}{1000^N},
 \qquad
 \alpha_{B,N}:=\operatorname{pref}_N(\alpha_B).
 \label{eq:alpha-definicion-b-y-prefijos-rev08}
\end{equation}

\begin{theorem}[Igualdad proyectiva de las dos publicaciones correlacionadas]
\label{prop:alpha-dos-limites}
Para todo \(N\), la vía entera y la vía analítica publican el mismo prefijo:
\begingroup
\renewcommand{\theequation}{A7\(_0\)}
\begin{equation}
 \boxed{\alpha_{A,N}=\operatorname{pref}_N(\alpha_A)
 =\operatorname{pref}_N(\alpha_B)=\alpha_{B,N}.}
 \label{eq:alpha-dos-vias-cada-nivel}
\end{equation}
\endgroup
En el límite,
\begingroup
\renewcommand{\theequation}{A8\(_0\)}
\begin{equation}
 \boxed{\alpha_A
 =\operatorname*{arg\,zero}_{x\in I_\alpha}
       E_{\mathsf s,\infty}(x)
 =\alpha_B=\alpha_{\mathrm{HMT}}.}
 \label{eq:alpha-dos-vias-limite-coinductivo}
\end{equation}
\endgroup
\end{theorem}

\begin{proof}
La normalización entera publica exactamente \(a(\mathsf s)\). El
teorema~\ref{thm:alpha-raiz-completa-explicita} demuestra que la carta
analítica representa ese mismo punto como raíz simple única. Los
cilindros semiabiertos \(C_N^A\) fijan un prefijo canónico y
\(I_{B,N}\subseteq C_N^A\) contiene \(\alpha_B=\alpha_A\). Por tanto, sus
prefijos coinciden para cada \(N\); el anidamiento y los diámetros tendentes a
cero dan la igualdad de los límites. El reconocimiento
metrológico ocurre después.
\end{proof}

La función analítica constituye una publicación correlacionada del mismo
estado, no una segunda derivación causalmente independiente de la clausura
entera. Esta precisión no reduce la igualdad demostrada: identifica su
mecanismo exacto, impide atribuir al jet una cola que no determina y evita la
promoción de un valor convencional a generador.

\endinput

% El bloque que sigue se conserva únicamente como antecedente editorial de
% REV04. \endinput impide que forme parte de la fuente activa o del PDF REV08.
\subsubsection{Lector completo, naturalidad y cilindros comunes}

Sea \(F_{\mathsf s}\) el cuerpo de coeficientes de la firma de
\(\mathsf s\in\mathcal X_\alpha^{\mathrm{enr}}\) y sea \(d_m=6+3m\). Se
definen
\[
 \mathfrak J_m:=\coprod_{\mathsf s\in\mathcal X_\alpha^{\mathrm{enr}}}
 \{\mathsf s\}\times F_{\mathsf s}[X]/(X^{d_m+1}),\qquad
 \Gamma_m(\mathsf s)=
 \bigl(\mathsf s,E_{21/22}^{(d_m)}(\mathsf s;X)\bigr).
\]
El lector completo registrado \(\Gamma_{E21}\) produce estas publicaciones.
Para \(m'\ge m\), el truncamiento
\(
 \tau_{m',m}(\mathsf s,[f]_{d_{m'}})
 =(\mathsf s,[f]_{d_m})
\)
satisface
\begin{equation}
 \tau_{m',m}\circ\Gamma_{m'}=\Gamma_m,
 \qquad
 E_{21/22}:=\varprojlim_mE_{21/22}^{(d_m)}.
 \tag{A3}
 \label{eq:alpha-sistema-analitico-completo}
\end{equation}
No se añade una sección del truncamiento desnudo. Si se explicita la
publicación de los tres coeficientes siguientes, ésta es la proyección
derivada
\[
 \Lambda_{d_m+1:d_m+3}
 :=\operatorname{pr}_{d_m+1:d_m+3}\circ
   \Gamma_{E21}\circ\mathcal U_{\mathrm{TPK}},
\]
no un generador independiente ni una lectura del valor objetivo.

Para un conjunto cofinal \(\mathcal N_\alpha\), el lector publica una función
cofinal \(m:\mathcal N_\alpha\to\mathbb N\) e intervalos
\(I_{B,N}=[\ell_N,u_N]\Subset I_\alpha\) tales que
\begin{equation}
\begin{gathered}
 I_{B,N'}\subseteq I_{B,N}\ (N'\ge N),\qquad
 \operatorname{diam}I_{B,N}\longrightarrow0,\qquad
 I_{B,N}\subseteq C_N^A,\\
 E_{21/22}^{(d_{m(N)})}(\ell_N)\le0\le
 E_{21/22}^{(d_{m(N)})}(u_N),\qquad
 \inf_{x\in I_{B,N}}\partial_xE_{21/22}^{(d_{m(N)})}(x)>0,\\
 E_{21/22}^{(d_{m(N)})}\longrightarrow E_{21/22}
 \quad\text{uniformemente en }I_\alpha .
\end{gathered}
 \tag{A4}
 \label{eq:alpha-via-b-proyectiva}
\end{equation}

\begin{proposition}[Cuadrados de compatibilidad y cilindros comunes]
\label{prop:alpha-cuadrados-compatibilidad}
Las restricciones de palabras y jets conmutan separadamente:
\[
\begin{array}{ccc}
 \mathcal X_\alpha^{\mathrm{enr}}&
 \xrightarrow{\ \mathsf T_{\alpha,N'}\ }&\mathcal W_{N'}\\[1mm]
 \Big\Vert&&\Big\downarrow\rho_{N',N}\\[1mm]
 \mathcal X_\alpha^{\mathrm{enr}}&
 \xrightarrow{\ \mathsf T_{\alpha,N}\ }&\mathcal W_N
\end{array}
\qquad
\begin{array}{ccc}
 \mathcal X_\alpha^{\mathrm{enr}}&
 \xrightarrow{\ \Gamma_{m'}\ }&\mathfrak J_{m'}\\[1mm]
 \Big\Vert&&\Big\downarrow\tau_{m',m}\\[1mm]
 \mathcal X_\alpha^{\mathrm{enr}}&
 \xrightarrow{\ \Gamma_m\ }&\mathfrak J_m .
\end{array}
 \tag{A5}
\]
La conexión entre ambos cuadrados es la inclusión certificada
\(I_{B,N}\subseteq C_N^A\), con anidamiento, contracción y cofinalidad. Por
ello, el punto único
\(x_B:=\bigcap_{N\in\mathcal N_\alpha}I_{B,N}\) satisface
\begin{equation}
 \forall N\in\mathcal N_\alpha,\qquad
 \operatorname{pref}_N(x_B)=\alpha_{A,N}=:\alpha_{B,N}.
 \tag{A6}
 \label{eq:alpha-igualdad-cada-nivel}
\end{equation}
\end{proposition}

\begin{proof}
La compatibilidad de la vía A es
\eqref{eq:alpha-via-a-proyectiva-exacta}; la de la vía B es
\eqref{eq:alpha-sistema-analitico-completo}. Cada jet posee una raíz simple
en su intervalo por el cambio de signo y la cota derivativa de
\eqref{eq:alpha-via-b-proyectiva}. La familia anidada determina \(x_B\), y la
convergencia uniforme da \(E_{21/22}(x_B)=0\). La inclusión cofinal en los
cilindros de la vía A fuerza sus prefijos. No se ha supuesto que truncar un
polinomio preserve su raíz.
\end{proof}

\begin{theorem}[Igualdad proyectiva y coinductiva de las dos vías]
\label{prop:alpha-dos-limites-historico-rev04}
Para todo nivel compatible del sistema cofinal común,
\begingroup
\renewcommand{\theequation}{A7}
\begin{equation}
 \boxed{\alpha_{A,N}=\rho_N(\alpha_A)
 =\rho_N(\alpha_B)=\alpha_{B,N}.}
 \label{eq:alpha-dos-vias-cada-nivel-historico-rev04}
\end{equation}
\endgroup
En el límite coinductivo,
\begingroup
\renewcommand{\theequation}{A8}
\begin{equation}
 \boxed{\alpha_A
 =\varprojlim_N\alpha_{A,N}
 =\varprojlim_N\alpha_{B,N}
 =\operatorname*{arg\,zero}_{x\in I_\alpha}E_{21/22}(x)
 =\alpha_B=\alpha_{\mathrm{HMT}}.}
 \label{eq:alpha-dos-vias-limite-coinductivo-historico-rev04}
\end{equation}
\endgroup
Ningún valor convencional interviene: el reconocimiento metrológico es
posterior a las dos publicaciones internas.
\end{theorem}

\begin{proof}
La vía A forma cilindros compatibles de diámetro \(1000^{-N}\); la vía B,
intervalos anidados de diámetro tendente a cero. Por
\eqref{eq:alpha-via-b-proyectiva}, el punto de la segunda familia es la raíz
simple y única del núcleo completo. Como pertenece a una subfamilia cofinal
de los cilindros de la primera, ambas intersecciones contienen el mismo punto.
La ecuación~\eqref{eq:alpha-igualdad-cada-nivel} da la igualdad a toda
profundidad compatible y el límite inverso da (A8). Las ejecuciones finitas
sólo certifican proyecciones; no acotan el enunciado.
\end{proof}

La raíz simple \(\xi_9\) y el testigo
\eqref{eq:alpha-testigo-finito-nueve} son verificaciones finitas de este
sistema. No definen el límite ni reducen (A7)--(A8) a una coincidencia
truncada: la igualdad completa procede de las extensiones del estado, de los
dos cuadrados de compatibilidad y de la unicidad de la intersección cofinal.

